\chapter*{Реферат}
\addcontentsline{toc}{chapter}{Реферат}

\textbf{Объект исследования} --- процесс автоматизированного планирования пешеходных
маршрутов с учётом индивидуальных предпочтений пользователя.

\textbf{Предмет исследования} --- математические методы оценки точек интереса,
построения взвешенного графа и поиска оптимального пути, а также программные
средства их реализации в веб-приложении.

\textbf{Цель работы} --- разработка прототипа веб-приложения Fun-Walk, способного
формировать маршруты прогулок по заданным критериям качества среды и визуализировать
результаты на интерактивной карте.

\textbf{Методы исследования:} системный анализ предметной области; линейная алгебра
(скалярное произведение векторов признаков и предпочтений); геодезические расчёты
на сфере (формула Haversine); теория графов и алгоритм Dijkstra; объектно-ориентированное
и модульное проектирование; клиент-серверная архитектура REST API.

\textbf{Результаты:} реализовано веб-приложение, включающее backend на NestJS и
frontend на React; разработана двухэтапная схема планирования (абстрактный граф
$\to$ OSRM-маршрутизация); проведено функциональное тестирование на демо-данных
центра Москвы.

\textbf{Объём работы:} \pageref{LastPage}~с., 22~источника, 1~приложение, 3~рисунка, 8~таблиц.

\textbf{Ключевые слова:} пешеходная маршрутизация, точки интереса, алгоритм Dijkstra,
Haversine, OSRM, NestJS, React, веб-приложение.

\clearpage
